% JAES Engineering Report
% aes2e documentclass — 2025 template
% Compiled with: pdflatex main.tex (x2)

\documentclass{aes2e}
\usepackage{capt-of}
% ── Required packages ───────────────────────────────────────────────
\usepackage{amsmath,amssymb}
\usepackage{booktabs}       % proper table rules
\usepackage{graphicx}
\usepackage{microtype}
\usepackage[section]{placeins}  % \FloatBarrier; restrict floats to their section
% hyperref must be loaded LAST; explicit color options prevent the
% [greenN]/[redN] color-name artifact some classes produce when
% defaults are left unspecified.
\usepackage[
  colorlinks=false,    % no colored boxes/text on links
  hidelinks,           % suppresses link borders entirely
  pdfborder={0 0 0},   % belt-and-suspenders: no border boxes
  hypertexnames=false  % avoids name collisions in some bibliography styles
]{hyperref}
\urlstyle{same}        % \url uses normal text font, not monospace

% ── Float placement parameters (balanced; works with [tp] specifier) ─
\renewcommand{\topfraction}{0.85}      % up to 85% of page may be top floats
\renewcommand{\textfraction}{0.10}     % keep at least 10% text on mixed pages
\renewcommand{\floatpagefraction}{0.75}% float-only pages must be 75%+ full
\setcounter{topnumber}{3}              % up to 3 top floats per page
\setcounter{totalnumber}{4}            % up to 4 floats total per page

% ── Metadata ─────────────────────────────────────────────────────────
\jyear{2026}
\jmonth{}
\jvol{}
\jnum{}

\begin{document}

\markboth{AHN}{REAL-TIME ADAPTIVE PEQ FOR ROOM CORRECTION AND TONAL SHAPE TRACKING}

% ── Title ────────────────────────────────────────────────────────────
% Note: "new", "novel" are prohibited. Title ≤ 125 chars.
% Char count: 104 chars (with spaces)
\title{Real-Time Adaptive Parametric Equalization for Joint Room Response
  Compensation and Tonal Shape Tracking%
  \thanks{Manuscript received 2026. Corresponding author: Yongho Ahn
    (nanhosoft@soongsil.ac.kr).}}

% ── Authors ──────────────────────────────────────────────────────────
\authorgroup{
  \author{Yongho Ahn},
  \role{Student Member, AES}
  \email{(nanhosoft@soongsil.ac.kr)}
  \affil{Department of AI Software Engineering, Soongsil University, Seoul, Republic of Korea}
}

% ── Abstract (100–200 words) ─────────────────────────────────────────
\abstract{
This engineering report describes a real-time adaptive parametric
EQ~(PEQ) that jointly performs room-response compensation and
operator-specified tonal-shape tracking. The task is objective spectral
target reproduction under room coloration; no subjective listening
preference claim is made. The model uses a lightweight TCN encoder with
RoomResponseEncoder and UserPreferenceEncoder conditioning, followed by a
two-stage head: Stage~A estimates a 128-bin room-correction prior and
Stage~B outputs 7-band PEQ parameters rendered through a differentiable
Gaussian approximation of peaking-EQ magnitude responses. Trained on
40{,}000 synthetic samples, the proposed model reaches LSD~1.442~dB,
DMR~0.928, and CosSim~0.960 on 3{,}000 held-out synthetic-RIR samples,
with inference-only RTF~$<$~0.001 on a CPU. Dense 128-bin architecture
variants achieve lower in-distribution LSD, but they do not emit PEQ
control parameters. The proposed 7-band output is deployment-compatible
and can be mapped to embedded PEQ control interfaces without fitting a
dense response. Real-RIR tests using BUT ReverbDB and OpenAIR show a
smaller synthetic-to-real gap for the proposed model than for the dense
variants and the auxiliary-loss variant.
}

\maketitle

% ═══════════════════════════════════════════════════════════════════════
\section{INTRODUCTION}
% ═══════════════════════════════════════════════════════════════════════

Music reproduction in everyday listening spaces faces two intertwined
challenges: room acoustics impose frequency-selective coloration on the
reproduced signal, and target reproduction may deviate from a flat
response according to operator-specified tonal shape requirements
(e.g., a bass-emphasis preset, vocal-clarity profile, or broadband
roll-off).

Commercial room-correction systems such as Dirac Live, Audyssey,
Sonos Trueplay, and Apple HomePod illustrate the practical importance of
combining room adaptation with user-facing tonal controls, although their
internal algorithms are generally not fully disclosed. In the hearing-aid
literature, preference-conditioned multi-band gain prediction has been
investigated extensively~\cite{keidser2006,nielsen2013,mondol2022};
these systems differ in target task (audibility compensation rather than
tonal shaping) and in output format (graphic-band gains rather than
biquad coefficients).

Within the published academic literature on neural parametric EQ, room
correction and preference-based EQ have largely been treated as separate
problems. Room-correction systems~\cite{miyoshi1988,pepe2020,pepe2022}
optimize for spectral flatness without operator-specified tonal targets;
preference- and text-based EQ systems~\cite{cartwright2013,seetharaman2016,
chu2025,stylianou2026} assume a flat reproduction environment.
Neural parametric-EQ parameter generation has been established in other
task settings~\cite{nercessian2020,nercessian2021_hc,colonel2022,
steinmetz2022}, and multi-task or conditioned formulations are a natural
next step that this report exercises end-to-end.

This report makes three practical contributions. First, it formulates a
dual-objective PEQ task that combines room compensation (LSD) with tonal
shape tracking (DMR, CosSim) in a single feed-forward model. Here, DMR
measures the fraction of frequency bins where the predicted response
matches the sign of the target tonal shape; it is an objective shape metric,
not a subjective preference score.

Second, it implements a lightweight \mbox{TCN$\times$4} encoder with
concatenated room and user-preference conditioning and a two-stage output
head: Stage~A estimates a 128-bin room-correction prior, and Stage~B emits
7-band PEQ parameters rendered with a differentiable Gaussian
approximation. These PEQ parameters ($f_c$, gain,~$Q$) can be mapped to
embedded PEQ control interfaces such as TAS5756M-style I$^2$C controls
without fitting a dense 128-bin response.

Third, it evaluates the system against six external baselines (E1--E6),
three dense-output architecture variants (AC1--AC3), targeted ablations
(A1, A2, A3), paired statistical diagnostics, track-level aggregate analysis
($N_\text{groups} = 1{,}306$), CPU RTF measurements, and real-RIR stress
tests using BUT ReverbDB~\cite{szoke2019} and OpenAIR~\cite{openair2010}.
A0 is not the highest-scoring architecture on in-distribution metrics:
AC variants achieve approximately 0.6~dB lower LSD with unconstrained
dense output. A0 is instead the deployment-compatible 7-band PEQ-output
configuration, and its measured synthetic-to-real gap is smaller than
those of the dense-output AC variants and the auxiliary-loss A2 variant.
The intent is to document a concrete implementation and evaluation
protocol for this engineering trade-off.

% ═══════════════════════════════════════════════════════════════════════
\section{RELATED WORK}
\label{sec:related}
% ═══════════════════════════════════════════════════════════════════════

\subsection{Room Acoustic Correction}

Classical room correction designs an inverse filter from a measured
room impulse response (RIR) to flatten the frequency
response~\cite{miyoshi1988}. Commercial products---Dirac Live,
Audyssey, Sonos Trueplay, Apple HomePod---show that microphone-assisted
room adaptation and user-facing target-curve controls are practically
important, but their internal algorithms are proprietary and not fully
auditable from public documentation.
In the academic literature, Pepe et al.~\cite{pepe2020} estimate an EQ
curve from audio features without explicit RIR input, and
BiasNet/CFN~\cite{pepe2022} from the same group couples differentiable
IIR parametric filters with a per-instance optimization targeting a flat
response. These systems optimize for spectral flatness and do not treat
an operator-specified tonal shape as an explicit model input.

\subsection{Personalized and Preference-Based Equalization}

Cartwright and Pardo's Social-EQ~\cite{cartwright2013} mapped
natural-language descriptors to EQ settings via crowdsourcing, later
extended to Audealize~\cite{seetharaman2016}. Recent text-driven EQ
leverages audio--language embeddings: Text2FX uses
CLAP~\cite{chu2025}, and Stylianou et al.~\cite{stylianou2026} use
LLM-based population-level tonal alignment.
The hearing-aid literature contains a substantial body of work on
preference-conditioned EQ personalization, including paired-comparison
learning of individual gain preferences~\cite{keidser2006,nielsen2013} and
deep-learning-based personalization of prescriptive fitting
rationales~\cite{mondol2022}; these systems output multi-band gain
vectors rather than parametric-biquad coefficients and target audibility
compensation rather than tonal shape reproduction in normal-hearing
listening.
These preference-oriented systems assume a flat or fixed reproduction
environment and do not perform room-response compensation jointly with the
preference mapping.

\subsection{Differentiable Parametric EQ}

Nercessian~\cite{nercessian2020} introduced a differentiable biquad
filter for neural PEQ parameter matching; the hyperconditioned
follow-up~\cite{nercessian2021_hc} mapped external user parameters to
biquad coefficients with stability-preserving activations for audio
distortion modelling. Colonel et al.~\cite{colonel2022} extended this
to direct biquad-cascade design from target magnitude responses;
Steinmetz et al.~\cite{steinmetz2022} (DeepAFx-ST) applied
differentiable DSP to audio-effect style transfer;
Sarkar and Lindborg~\cite{sarkar2025} (Diff-DEQ) combined
FiLM-modulated TCN with BiGRU for adaptive dynamic equalization;
and Mockenhaupt et al.~\cite{mockenhaupt2024} applied differentiable
PEQ to per-instrument automatic equalization. These studies establish
that neural PEQ parameter generation is feasible and deployable, but
each targets a different task than the joint
room-correction/tonal-shape setting here.

\subsection{Positioning of This Work}

Compared with the work surveyed above, the implementation described in
this report combines four properties: (i)~feed-forward inference that
generalizes across unseen rooms and preferences without per-instance
retraining, (ii)~an explicit user-preference input branch fused with a
room-response branch in a single jointly trained network, (iii)~a
PEQ-parameter output head rendered by a differentiable Gaussian approximation
whose parameters can be mapped to embedded PEQ control interfaces, and
(iv)~release of implementation scripts and generated evaluation artifacts
where licensing permits.
Individual elements of this combination appear in prior work and commercial
products; the engineering contribution of this report is the integration
into a single documented implementation and evaluation protocol.

% ═══════════════════════════════════════════════════════════════════════
\section{PROPOSED METHOD}
% ═══════════════════════════════════════════════════════════════════════

\subsection{Problem Formulation}

Let $\mathbf{x} \in \mathbb{R}^T$ denote a music signal reproduced in a
room with impulse response $h$.
The observed signal is $y = x * h$, where the room induces spectral
coloration $H(f) = Y(f)/X(f)$.
Simultaneously, the desired reproduction target deviates from a flat response
according to an operator-specified tonal shape $P(f)$ (e.g., a bass-emphasis
curve, a vocal-clarity profile, or a broadband roll-off).
$P(f)$ represents a prescribed spectral target, not a claim about
listener preference.

We formulate the dual-objective EQ problem as follows: given acoustic features
extracted from $y$ and a user tonal shape specification (mode
$m \in \{0,1,2,3\}$ and per-band gain vector $\mathbf{g} \in \mathbb{R}^{10}$),
predict an EQ response $R(f)$ such that the resulting response simultaneously
compensates room-induced spectral distortion and tracks the target tonal
shape $P(f)$ with high directional fidelity.
The dual target used for training and for the headline tables is:
\begin{equation}
  T_\text{dual}(f) = T_\text{room}(f) + T_\text{pref}(f),
\end{equation}
where $T_\text{room}(f)$ is the ideal room correction response and
$T_\text{pref}(f)$ is the tonal shape profile interpolated from the
band-gain vector $\mathbf{g}$ at 10 graphic-EQ center frequencies
[63, 125, 250, 500, 1k, 2k, 4k, 8k, 12k, 16k]~Hz.

\subsection{Model Architecture}

Fig.~\ref{fig:arch} shows the proposed DualObjectiveAdaptivePEQ architecture.
The model consists of six components.
\begin{figure}[tp]
  \centering
  \includegraphics[width=\linewidth]{fig1_architecture_docx.png}
  \caption{Two-stage DualObjectiveAdaptivePEQ architecture overview. The RoomResponseEncoder (orange), TCN$\times$4 backbone (green), and UserPreferenceEncoder (blue) feed Stage~A (room correction prior) and Stage~B (preference-conditioned PEQ).}
  \label{fig:arch}
\end{figure}


\textbf{RoomResponseEncoder.}
A two-layer MLP (128$\to$64$\to$32, LayerNorm$+$GELU) compresses the
128-bin room response input into a room vector
$\mathbf{v}_r \in \mathbb{R}^{32}$.
Ablations confirm that removing this component degrades
out-of-distribution robustness (removing it raises real-RIR LSD to 4.875~dB vs.\ proposed: 1.941~dB).

\textbf{UserPreferenceEncoder.}
A mode embedding ($n_\text{modes} = 4$, dim~$= 16$) is concatenated with a
projected band-gain vector (10$\to$64, MLP$+$LayerNorm$+$GELU), and fused
to produce $\mathbf{v}_p \in \mathbb{R}^{64}$.
Conditioning is applied by concatenation in the reported model rather than
via FiLM~\cite{perez2018}.

\textbf{Input Projection.}
Acoustic features $\mathbf{x} \in \mathbb{R}^{T \times 10}$ are projected
to $\mathbb{R}^{T \times 64}$ via a linear layer with LayerNorm and GELU.

\textbf{TCN Blocks.}
Four causal dilated temporal convolutional blocks
($n_\text{blocks} = 4$, dilation $\in \{1,2,4,8\}$, kernel size~$= 3$)
process the sequence with residual connections.
Left-padding only ensures no future leakage, maintaining compatibility
with segment-based real-time deployment.

\textbf{Attention Pooling.}
A learnable AttentionPooling module aggregates the sequence into context
vector $\mathbf{ctx} \in \mathbb{R}^{64}$.

\textbf{Two-Stage Output Head.}
Stage~A (room correction prior): $\mathbf{v}_r$ alone is fed to
\texttt{room\_mean\_head} (scalar mean-shift, scale~$= 4.0$~dB) and
\texttt{room\_shape\_head} (128-bin mean-centered shape,
scale~$= 12.0$~dB), which are summed to produce
$\mathbf{r}_\text{corr} \in \mathbb{R}^{128}$.
Stage~A deliberately excludes $\mathbf{ctx}$ to ensure that room correction
depends solely on the measured room response, independent of musical content.
This contrasts with the AC variants, which condition Stage~A on both
$\mathbf{ctx}$ and $\mathbf{v}_r$; the $\mathbf{v}_r$-only design in A0
is motivated by improving generalization to unseen acoustic
environments; this design is supported empirically by A0's
smaller synthetic-to-real domain gap (Section~\ref{sec:ood}).

Stage~B (preference-conditioned PEQ):
$[\mathbf{ctx},\, \mathbf{v}_r,\, \mathbf{v}_p,\, \mathbf{r}_\text{corr},\,
\mathbf{p}_\text{curve}]$ are fed to \texttt{peq\_head}, which outputs
$K=7$ filter parameters ($f_c \in [80, 16000]$~Hz, gain $\in [-6, +6]$~dB,
$Q \in [0.3, 8.0]$) per filter.
A differentiable Gaussian approximation of the peaking-EQ magnitude response
(following Nercessian~\cite{nercessian2020}) renders these into a 128-bin
frequency response $\mathbf{r}_\text{PEQ}$.
The final prediction is:
\begin{equation}
  \hat{R}(f) = \mathbf{r}_\text{corr}(f) + \mathbf{r}_\text{PEQ}(f).
\end{equation}
This design provides a deployment-compatible control output: the $K=7$
filter parameters can be mapped to embedded PEQ control interfaces
(e.g., TAS5756M-style I$^2$C controls) without fitting a dense 128-bin
response after inference. The \texttt{peq\_head} outputs raw parameters
that are sigmoid/tanh-mapped to physical ranges, then converted to a
128-bin magnitude response through the Gaussian PEQ approximation.
The two-stage factorization encourages the model to first address room
correction (Stage~A) and then apply parametric adjustments to track the
residual tonal shape target (Stage~B).

\subsection{Loss Function}

The proposed model is trained with a two-term top-level loss:
\begin{equation}
  \mathcal{L} = \lambda_\text{final}\,\mathcal{L}_\text{final} +
                \lambda_\text{room}\,\mathcal{L}_\text{room},
\end{equation}
where $\lambda_\text{final}=1.0$ and $\lambda_\text{room}=0.35$.
$\mathcal{L}_\text{final}$ is a weighted curve loss between
$\hat{R}$ and $T_\text{dual}$, incorporating A-weighting-based perceptual
weights, low/high-edge frequency boosts, shape/gradient/curvature
sub-terms, and a small mean-level penalty ($\lambda_\text{mean}=0.03$)
to discourage global level drift between the predicted and target
responses.
$\mathcal{L}_\text{room}$ directly supervises $\mathbf{r}_\text{corr}$
against $T_\text{room}$, encouraging Stage~A to specialize in room
correction.
A more elaborate variant adds explicit Stage-B supervision via a
preference-residual loss $\mathcal{L}_\text{pref}$ (on the PEQ residual
against $T_\text{dual} - T_\text{room}$) and a directional loss
$\mathcal{L}_\text{dir} = 1 - \text{CosSim}(\hat{R} - T_\text{room},
T_\text{pref})$.
This four-term variant was investigated and is reported as ablation
``A2 with Pref Loss'' in Section~\ref{sec:ablation}; it does not improve
test-set metrics over the simplified two-term top-level loss used in the
proposed model and is therefore not part of A0.

% ═══════════════════════════════════════════════════════════════════════
\section{EXPERIMENTAL SETUP}
% ═══════════════════════════════════════════════════════════════════════

\subsection{Dataset}

We constructed a synthetic dataset by convolving FMA-Medium music
clips~\cite{defferrard2017} (25{,}000 tracks, 16~genres) with simulated RIRs
generated using pyroomacoustics~\cite{scheibler2018}.
Room parameters were sampled uniformly: room size 3.5--8.0~m,
RT60 0.2--0.75~s, source--microphone distance 0.7--2.2~m.
These ranges reflect realistic domestic listening environments (living rooms,
offices).
Microphone estimation noise ($\sigma \in [0.4, 2.0]$~dB, bias $\leq 1.25$~dB)
was added to simulate imperfect room response measurements.

Four tonal shape presets (Vocal, Bass, Treble, Soft) were manually defined to
span representative spectral emphasis patterns across the 63~Hz--16~kHz range:
low-frequency emphasis~(Bass), upper-midrange boost~(Vocal),
high-frequency emphasis~(Treble), and broadband roll-off~(Soft).
These presets serve as illustrative examples of operator-specified tonal shapes;
the model and training procedure are independent of the specific preset
definitions, and any user-defined target tonal shape can be substituted.
During target generation, the per-sample band-gain input is folded into the
synthesized preference target with a partial weight of 0.55, modelling soft
control compliance rather than verbatim copying of the operator-provided
band vector. The preset center-band values used as the base tonal profiles are
listed in Table~\ref{tab:profiles}.

\begin{table}[tp]
  \centering
  \caption{Tonal shape preset center-band gains (dB).
           Frequency bands in Hz.}
  \label{tab:profiles}
  \resizebox{\linewidth}{!}{%
  \begin{tabular}{lrrrrrrrrrr}
    \toprule
    Mode & 63 & 125 & 250 & 500 & 1k & 2k & 4k & 8k & 12k & 16k \\
    \midrule
    Vocal  & $-4$ & $-3$ & $0$ & $+2$ & $+5$ & $+6$ & $+4$ & $+2$ & $-1$ & $-2$ \\
    Bass   & $+8$ & $+6$ & $+3$ & $0$ & $-1$ & $-2$ & $-3$ & $-2$ & $-1$ & $0$ \\
    Treble & $-2$ & $-1$ & $0$ & $0$ & $+1$ & $+3$ & $+5$ & $+7$ & $+8$ & $+6$ \\
    Soft   & $+2$ & $+1$ & $+1$ & $0$ & $-1$ & $-3$ & $-5$ & $-6$ & $-4$ & $-3$ \\
    \bottomrule
  \end{tabular}%
  }
\end{table}

Data were split by track identity into train, validation, and held-out
synthetic test partitions to prevent within-track leakage across splits
(40{,}000~/~5{,}000~/~3{,}000 samples respectively).
Each sample consists of a 4-second clip convolved with an independently
sampled RIR; multiple clips may be drawn from the same track but never
appear in more than one split.
All experiments used a fixed random seed (seed~$= 42$).
An additional out-of-distribution test set of 2{,}000 samples was constructed
using real-world RIRs from BUT ReverbDB~\cite{szoke2019} and OpenAIR~\cite{openair2010}.
All audio was resampled to 48~kHz.
The frequency-response targets are represented as 128 log-spaced bins
between 20~Hz and 24~kHz; the same grid is shared across the dataset
generator, model output, baselines, and loss computation.
The synthetic dataset generation pipeline, model code, and training scripts
are implemented in PyTorch; code availability will be stated upon acceptance.

\subsection{Input Features}

From each 4-second clip we extract a $T = 32$ frame $\times$ 10-dimensional
feature sequence: 8~MFCC coefficients, 1~log-energy, and 1~normalized spectral
centroid.
Features are extracted with a 25~ms window and 10~ms hop
(\texttt{hop\_length} $= 480$ samples at 48~kHz), then reduced to
32~frames by uniform temporal index sampling, followed by per-dimension
$z$-score normalization.
The 32-frame resolution is sufficient for capturing room-level spectral
characteristics, which vary slowly relative to musical transients;
fine-grained temporal detail is provided directly by the 128-bin
\texttt{room\_response} input.

\subsection{Evaluation Metrics}

We report three variants of LSD:
$\text{LSD}_\text{dual}$ (predicted EQ response vs.\ $T_\text{dual}$,
primary metric), $\text{LSD}_\text{room}$ (predicted EQ response vs.\
$T_\text{room}$), and $\text{LSD}_\text{pref}$ (heard response vs.\
$T_\text{pref}$). The estimated heard response is defined as
$\hat{H}(f) = \hat{R}(f) - T_\text{room}(f)$, i.e., the predicted EQ
response after accounting for the room-correction target. Unless otherwise
noted, LSD refers to $\text{LSD}_\text{dual}$.

\textbf{Directional Match Rate~(DMR)} is the fraction of frequency bins where
the heard response after room compensation has the same sign as the tonal
shape target:
$\text{sign}(\hat{R}(f)-T_\text{room}(f)) =
\text{sign}(T_\text{pref}(f))$, higher is better.
DMR directly measures tonal shape tracking fidelity---the degree to which
the predicted response reproduces the directional structure of the target
shape at each frequency bin after accounting for room correction.
It is a purely objective metric and does not imply subjective preference
alignment.
It is complemented by \textbf{Cosine Similarity~(CosSim)}, computed between
the heard response and the target tonal shape.
\textbf{Real-Time Factor~(RTF$^\dagger$)} $=$ inference-only forward-pass
latency divided by the 4-second segment duration used in the evaluation
script; RTF $< 1.0$ indicates real-time inference capability.
($^\dagger$ Excludes the time needed to accumulate the 4-second audio
context and to extract acoustic features and the room-response input;
see Section~\ref{sec:rtf}.)

\subsection{Baselines and Comparisons}

Six external baselines span the design space.
E1~(No Processing) serves as a lower bound.
E2~(Static Mode EQ) applies a fixed rule-based profile per mode without
room awareness and does not use the per-sample \texttt{band\_gains} vector.
E3~(Nercessian MLP~\cite{nercessian2020}) and E4~(Pepe CNN~\cite{pepe2020})
are reimplementations of prior neural room correction systems, trained on
$T_\text{room}$ only.
The primary limitation of E3 and E4 as baselines is not feature dimensionality
but objective mismatch: both models are trained on room correction targets
only, whereas A0 optimizes a dual target $T_\text{dual}$.
This mismatch is intentional---it reflects the original design scope of
these systems---and explains their near-random DMR scores regardless of
architecture.

E5~(Sequential: E3$\to$E2) represents the na\"{i}ve combination of room
correction and tonal shape preset EQ.
E6~(DSP Analytical) inverts the measured room response and adds the tonal
shape profile analytically, with no learned parameters; like E2, it uses
the discrete mode profile but not the per-sample \texttt{band\_gains}
conditioning.
E2 and E6 are therefore interpreted as intentionally simple static or
analytical reference operating points rather than capacity-matched
full-conditioning competitors to A0. Conditioning-aware learned comparisons
are provided by A3 and AC1--AC3.


Three architecture comparison variants (AC1--AC3) share the same
dual-objective loss and concat conditioning as the proposed model, but differ
from A0 in three respects: (i)~the sequence aggregation mechanism,
(ii)~Stage~A conditioning (\texttt{ctx}$+$\texttt{room\_vec}
vs.\ \texttt{room\_vec} only in A0), and (iii)~Stage~B output structure
(direct 128-bin via \texttt{pref\_res\_head}, rather than A0's
PEQ-parameterized Stage~B).
AC1~(TCN$+$BiLSTM) uses bidirectional LSTM mean pooling;
AC2~(TCN$+$GRU) uses causal single-direction GRU;
AC3~(TCN$+$Conformer~\cite{gulati2020}) uses two Conformer blocks followed
by attention pooling.

\subsection{Ablation Design}

Three targeted ablations isolate the contribution of each key design decision.
A1~(w/o Room Awareness) removes both the explicit \texttt{room\_response}
input and the room-induced coloration cues in the audio feature stream: when
clean pre-convolution features are available, A1 uses them instead of the
reverberant feature sequence. This ablation therefore tests dependence on
room-related information as a whole---explicit room measurements and implicit
reverberation cues---rather than only the explicit room-response vector.
A2~(+Pref/Dir Loss) keeps the same two-stage PEQ-parameter architecture as A0
but adds explicit Stage-B preference residual supervision
$\mathcal{L}_\text{pref}$ and directional alignment loss
$\mathcal{L}_\text{dir}$. This ablation tests whether auxiliary
preference-specific losses improve over direct supervision of the final
dual target.
A3~(w/o Pref Input) removes both \texttt{mode\_id} and \texttt{band\_gains}
inputs entirely, eliminating the UserPreferenceEncoder conditioning---this
isolates whether tonal shape tracking originates from input conditioning or
from loss design.

All trainable models share the same training protocol: AdamW
optimizer~($lr = 3 \times 10^{-4}$, weight\_decay~$= 10^{-4}$),
CosineAnnealingLR~($T_\text{max} = 300$, $\eta_\text{min} = 10^{-6}$),
batch size~512, early stopping with patience~$= 15$ epochs
monitored on val\_lsd.

% ═══════════════════════════════════════════════════════════════════════
\section{RESULTS}
% ═══════════════════════════════════════════════════════════════════════

\subsection{Main Results}

Table~\ref{tab:main} summarizes results across all models on a held-out
test set of 3{,}000 synthetic-RIR samples.
The two neural baselines that target only room correction (E3, E4)
record near-random DMR (0.507 and 0.504), demonstrating that room
correction alone systematically fails to track the target tonal shape.
The static or sequential baselines that include a preset-based EQ
component (E2, E5) reach DMR 0.736 and 0.767 but pay a substantial LSD
penalty (3.670--3.573~dB) under room-degraded conditions.
The analytical baseline E6, which inverts the measured room response
and adds the tonal preset analytically, is the strongest non-learned baseline
(LSD~2.330~dB, DMR~0.882).
The three architectural variants AC1--AC3, which share the proposed
encoder backbone but use a dense 128-bin output head, obtain the
lowest in-distribution LSD (0.853--0.870~dB) and the highest DMR
(0.945--0.947).
The proposed model A0 achieves LSD~1.442~dB, DMR~0.928, CosSim~0.960,
trailing the AC variants by approximately 0.6~dB on dual-target LSD.
Unlike the AC variants, A0 emits 7-band PEQ parameters~($f_c$, gain,
$Q$) that can be mapped to a target amplifier control interface without
post-hoc dense-response fitting.
A0 also remains closer to the room-only target under the Room-LSD
diagnostic (LSD$_\text{room}$~4.689~dB) than the AC variants
(5.226--5.238~dB), consistent with its \texttt{room\_vec}-only
factorization.
All evaluated models remain at or below RTF~0.001 under the
inference-only 4-second-segment definition.


Fig.~\ref{fig:room_correction} examines room correction behaviour as a function of room coloration severity and reverberation time. Fig.~\ref{fig:per_mode} breaks results down per tonal shape preset, showing that A0 maintains consistent performance across all four modes while baselines collapse to near-random DMR.
\begin{figure}[tp]
  \centering
  \includegraphics[width=\linewidth]{fig2_room_correction.pdf}
  \caption{Room correction performance: (a) mean predicted response vs.\ room coloration bias; (b) LSD by RT60 quartile.}
  \label{fig:room_correction}
\end{figure}

\begin{figure}[tp]
  \centering
  \includegraphics[width=\linewidth]{fig3_per_mode.pdf}
  \caption{Per-mode LSD (top) and DMR (bottom) for E3, E4, and A0.}
  \label{fig:per_mode}
\end{figure}


\begin{table}[tp]
  \centering
  \caption{Main results on the held-out test set of 3{,}000 synthetic-RIR
           samples. \textbf{Bold} denotes the best result per metric.
           RTF$^\dagger$ = evaluation-script inference-only latency /
           4-second segment duration (excludes context accumulation; see
           Section~\ref{sec:rtf}). All models measured on AMD Ryzen
           9800X3D~(single-threaded). $\star$~denotes the proposed model.}
  \label{tab:main}
  \resizebox{\linewidth}{!}{%
  \begin{tabular}{lccccc}
    \toprule
    Model & LSD $\downarrow$ (dB) & DMR $\uparrow$ & CosSim $\uparrow$
          & RTF$^\dagger$ $\downarrow$ & Params \\
    \midrule
    E1~No Processing            & 6.033 & 0.482 & 0.027 & $<$0.001 & 0 \\
    E2~Static Mode EQ           & 3.670 & 0.736 & 0.720 & $<$0.001 & 0 \\
    E3~Nercessian MLP~\cite{nercessian2020} & 5.964 & 0.507 & 0.011 & $<$0.001 & 138{,}255 \\
    E4~Pepe CNN~\cite{pepe2020}             & 5.809 & 0.504 & 0.032 & $<$0.001 & 50{,}351 \\
    E5~Sequential (E3$\to$E2)   & 3.573 & 0.767 & 0.732 & $<$0.001 & 138{,}255 \\
    E6~DSP Analytical           & 2.330 & 0.882 & 0.919 & $<$0.001 & 0 \\
    \midrule
    AC1~TCN$+$BiLSTM            & \textbf{0.853} & \textbf{0.947} & \textbf{0.982} & 0.001    & 247{,}713 \\
    AC2~TCN$+$GRU               & 0.860 & 0.946 & 0.982 & $<$0.001 & 272{,}545 \\
    AC3~TCN$+$Conformer~\cite{gulati2020} & 0.870 & 0.945 & 0.982 & $<$0.001 & 418{,}914 \\
    \midrule
    A0~Proposed $\star$         & 1.442 & 0.928 & 0.960 & $<$0.001 & 203{,}447 \\
    A1~w/o Room Awareness           & 2.897 & 0.832 & 0.862 & $<$0.001 & 140{,}033 \\
    A2~with Pref Loss           & 1.703 & 0.884 & 0.944 & $<$0.001 & 203{,}447 \\
    A3~w/o Pref Input           & 5.329 & 0.509 & 0.084 & $<$0.001 & 140{,}257 \\
    \bottomrule
  \end{tabular}%
  }
\end{table}

\subsection{Ablation Study}
\label{sec:ablation}

Table~\ref{tab:ablation} reports ablation results.
Ablation A3 (w/o Pref Input) produces the largest degradation overall:
LSD rises from 1.442 to 5.329~dB and DMR drops from 0.928 to 0.509,
with paired effect size $d_z = +2.376$ in favour of A0 on DMR
($p < 0.001$).
This identifies the UserPreferenceEncoder input as the primary driver of
tonal shape tracking fidelity.
Ablation A1 (w/o Room Awareness) yields LSD~2.897~dB and DMR~0.832 on
synthetic RIRs, then degrades further to LSD~4.875~dB and DMR~0.740 on
real RIRs (versus proposed A0: 1.941~dB, 0.885), confirming that
room-related information---explicit room measurements and implicit
reverberation cues---is essential both for in-distribution accuracy and
for out-of-distribution robustness.
Ablation A2 (with Pref Loss) is identical to A0 in architecture but
adds the explicit Stage-B preference loss $\mathcal{L}_\text{pref}$ and
direction loss $\mathcal{L}_\text{dir}$ described in Section~3.3.
A2 yields LSD~1.703~dB and DMR~0.884, both worse than the simplified
A0 (LSD~1.442~dB, DMR~0.928); the proposed model attains 0.26~dB
lower LSD and 0.044 higher DMR than A2 ($d_z = -0.486$ for LSD,
$d_z = +0.404$ for DMR; both $p < 0.001$).
The small but consistent advantage of A0 over A2 indicates that the
explicit Stage-B loss terms do not improve test-time metrics for this
task; tonal shape tracking fidelity emerges primarily from the input
conditioning structure (UserPreferenceEncoder) and the additive
dual-target formulation $T_\text{dual}$ rather than from auxiliary loss
terms.


Fig.~\ref{fig:ablation_bar} summarizes the ablation metric shifts relative to the A0 baseline.
\begin{figure}[tp]
  \centering
  \includegraphics[width=\linewidth]{fig4_ablation_bar.pdf}
  \caption{Ablation study: LSD and DMR per variant. Dashed line = A0 baseline.}
  \label{fig:ablation_bar}
\end{figure}


\begin{table}[tp]
  \centering
  \caption{Ablation results on the test set ($N=3{,}000$).
           $^\ddagger$$p < 0.001$ vs.\ proposed A0 in raw paired tests.}
  \label{tab:ablation}
  \resizebox{\linewidth}{!}{%
  \begin{tabular}{llccc}
    \toprule
    Variant & Description & LSD $\downarrow$ & DMR $\uparrow$ & CosSim $\uparrow$ \\
    \midrule
    A0~$\star$ & Proposed (simplified two-term loss)
               & \textbf{1.442} & \textbf{0.928} & \textbf{0.960} \\
    A1         & w/o Room Awareness
               & 2.897 & 0.832$^\ddagger$ & 0.862$^\ddagger$ \\
    A2         & with Pref Loss (adds $\mathcal{L}_\text{pref}+\mathcal{L}_\text{dir}$)
               & 1.703 & 0.884$^\ddagger$ & 0.944$^\ddagger$ \\
    A3         & w/o Pref Input
               & 5.329 & 0.509$^\ddagger$ & 0.084$^\ddagger$ \\
    \bottomrule
  \end{tabular}%
  }
\end{table}

\subsection{Paired Statistical Diagnostics}
\label{sec:stats}

Table~\ref{tab:stats} reports raw paired-test diagnostics on the
held-out test set ($N=3{,}000$).
These p-values are not multiple-comparison adjusted; LSD is treated as the
primary error metric, while DMR is reported as a secondary tonal-alignment
diagnostic.
All comparisons against the room-only neural baselines (E3, E4) show
small raw paired-test $p$-values ($p < 0.001$) with large effect sizes on both metrics:
A0 attains 4.522~dB lower LSD than E3 ($d_z = -3.702$) and 4.366~dB
lower than E4 ($d_z = -4.033$); A0's DMR exceeds E3 by 0.422 and E4 by
0.425 ($d_z = +2.329$ and $+2.431$, both $p < 0.001$).
A0 vs.\ A3 (w/o Pref Input) yields the largest overall effect size
(LSD: 3.886~dB lower, $d_z = -4.226$; DMR: 0.420 higher,
$d_z = +2.376$), indicating that the UserPreferenceEncoder input is the
single most impactful component in these diagnostics.
A0 vs.\ A1 (w/o Room Awareness) also shows a large raw paired-test difference
(LSD: 1.455~dB lower, $d_z = -1.082$; DMR: 0.096 higher,
$d_z = +0.806$, both raw $p < 0.001$).
A0 vs.\ A2 (with Pref Loss) also yields small raw paired-test $p$-values
($p < 0.001$) but with small effect sizes: A0 attains 0.261~dB lower LSD
($d_z = -0.486$) and 0.044 higher DMR ($d_z = +0.404$) than A2,
consistent with the interpretation that the explicit Stage-B loss
terms provide marginal, and on this test set slightly negative,
contributions to inference metrics.
A0 vs.\ AC variants also yields small raw paired-test $p$-values
($p < 0.001$) but with small effect sizes in the opposite direction: AC1--AC3 attain 0.57--0.59~dB
lower LSD ($d_z \approx +1.16$ to $+1.19$ in favour of AC) and 0.017--0.019
higher DMR ($d_z \approx -0.54$ to $-0.56$); these reflect the
in-distribution advantage of unconstrained dense 128-bin output over
the proposed 7-band PEQ parameter output, discussed further in
Section~\ref{sec:arch}.

Track-level aggregate analysis ($N_\text{groups} = 1{,}306$ tracks)
addresses potential within-track correlation by collapsing each
track's samples to a single per-track mean before pairing.
The findings are consistent with the sample-level pattern.
On LSD, the proposed A0 achieves a 1.254~dB lower per-track mean than
A1 (Cohen's $d_z$ in the direction favouring A0; track-level
raw $p < 0.001$), 0.257~dB lower than A2,
3.890~dB lower than A3, 4.542~dB lower than E3, and 4.380~dB lower
than E4.
On DMR, A0 attains 0.084 higher per-track mean than A1, 0.044 higher
than A2, 0.422 higher than A3, 0.427 higher than E3, and 0.428 higher
than E4 (all raw paired-test $p < 0.001$).
The track-level effect sizes corroborate the sample-level pattern:
A0 outperforms the room-only or static-EQ baselines and ablations in these
paired diagnostics, while the dense-output AC variants attain 0.574--0.592~dB
lower per-track LSD than A0 and 0.016--0.018 higher per-track DMR.

Fig.~\ref{fig:dmr_scatter} shows the per-sample DMR gap between A0 and
two informative contrasts (A3, baseline E3); Fig.~\ref{fig:per_mode_response}
overlays predicted frequency responses by tonal shape preset across 200
test samples.
\begin{figure}[tp]
  \centering
  \includegraphics[width=\linewidth]{fig6_scatter.pdf}
  \caption{Sample-level DMR scatter: A0 vs.\ A3 (w/o Pref Input) and A0 vs.\ E3.}
  \label{fig:dmr_scatter}
\end{figure}

\begin{figure}[tp]
  \centering
  \includegraphics[width=\linewidth]{fig7_freq_response.pdf}
  \caption{Per-mode predicted frequency responses (200 test samples,
           shaded region = $\pm 1$ standard deviation).}
  \label{fig:per_mode_response}
\end{figure}


\begin{table}[tp]
  \centering
  \caption{Paired statistical test results on the test set
           ($N = 3{,}000$). Raw paired-test $p$-values are reported as diagnostic statistics.
           $\Delta$ favours A0 sign convention: LSD shown as
           A0 minus baseline (negative = A0 lower LSD = A0 better);
           DMR shown as A0 minus baseline (positive = A0 higher DMR
           = A0 better).
           $d_z$ = Cohen's $d_z$ for paired comparisons.
           Win\% = fraction of test samples where A0 is strictly better.}
  \label{tab:stats}
  \resizebox{\linewidth}{!}{%
  \begin{tabular}{llcccc}
    \toprule
    Comparison & Metric & $\Delta$ mean & $p$-value & $d_z$ & Win\% \\
    \midrule
    A0 vs.\ A1 & LSD & $-1.455$ & $<$0.001 & $-1.082$ & 86.5\% \\
    A0 vs.\ A1 & DMR & $+0.096$ & $<$0.001 & $+0.806$ & 78.3\% \\
    \midrule
    A0 vs.\ A2 & LSD & $-0.261$ & $<$0.001 & $-0.486$ & 67.3\% \\
    A0 vs.\ A2 & DMR & $+0.044$ & $<$0.001 & $+0.404$ & 49.6\% \\
    \midrule
    A0 vs.\ A3 & LSD & $-3.886$ & $<$0.001 & $-4.226$ & 100.0\% \\
    A0 vs.\ A3 & DMR & $+0.420$ & $<$0.001 & $+2.376$ & 99.5\% \\
    \midrule
    A0 vs.\ E3 & LSD & $-4.522$ & $<$0.001 & $-3.702$ & 99.7\% \\
    A0 vs.\ E3 & DMR & $+0.422$ & $<$0.001 & $+2.329$ & 99.1\% \\
    \midrule
    A0 vs.\ E4 & LSD & $-4.366$ & $<$0.001 & $-4.033$ & 100.0\% \\
    A0 vs.\ E4 & DMR & $+0.425$ & $<$0.001 & $+2.431$ & 99.3\% \\
    \midrule
    A0 vs.\ AC1 & LSD & $+0.589$ & $<$0.001 & $+1.194$ & 7.6\% \\
    A0 vs.\ AC1 & DMR & $-0.019$ & $<$0.001 & $-0.560$ & 22.8\% \\
    A0 vs.\ AC2 & LSD & $+0.582$ & $<$0.001 & $+1.175$ & 8.3\% \\
    A0 vs.\ AC2 & DMR & $-0.018$ & $<$0.001 & $-0.550$ & 23.4\% \\
    A0 vs.\ AC3 & LSD & $+0.572$ & $<$0.001 & $+1.159$ & 8.7\% \\
    A0 vs.\ AC3 & DMR & $-0.017$ & $<$0.001 & $-0.539$ & 24.2\% \\
    \bottomrule
  \end{tabular}%
  }
\end{table}

\subsection{Out-of-Distribution Generalization}
\label{sec:ood}

Table~\ref{tab:ood} reports out-of-distribution generalization results
using real-world RIRs from BUT ReverbDB~\cite{szoke2019} and OpenAIR~\cite{openair2010} ($N = 2{,}000$).
All preference-aware models exhibit modest LSD increases on real RIRs.
The proposed A0 shows a smaller synthetic-to-real domain gap than the
dense-output AC variants and the auxiliary-loss A2 variant (0.499~dB),
reaching real LSD~1.941~dB and real DMR~0.885; A2 (with Pref Loss)
follows with a 0.540~dB gap, and the AC variants show 0.646--0.706~dB
gaps despite their lower in-distribution LSD.
The critical finding is A1 (w/o Room Awareness): when both the explicit
room-response input and implicit reverberation cues in the feature stream
are removed, real-RIR LSD reaches 4.875~dB, 2.5$\times$ higher than A0's,
with a much larger domain gap (1.978~dB). This confirms that room-awareness
information is the primary mechanism for acoustic generalization.
This smaller gap relative to the AC and A2 variants indicates that the
\texttt{room\_vec}-only Stage~A conditioning, which excludes the audio
\texttt{ctx} from the room-correction prior, limits synthetic-to-real
drift under acoustic environments unseen during training.
Many real-RIR rooms in BUT ReverbDB and OpenAIR exceed the training
distribution (RT60~0.20--0.75~s), making this a deliberate stress test
rather than an in-distribution evaluation.

Table~\ref{tab:ood} summarizes the LSD and DMR breakdown under the two
test conditions.




\begin{table}[tp]
  \centering
  \caption{Out-of-distribution generalization: held-out synthetic-RIR
           test set ($N=3{,}000$) vs.\ real-world RIRs from
           BUT ReverbDB~\cite{szoke2019} and OpenAIR~\cite{openair2010} ($N = 2{,}000$).
           \textbf{Bold} denotes best result per column.}
  \label{tab:ood}
  \resizebox{\linewidth}{!}{%
  \begin{tabular}{lcccc}
    \toprule
    Model & Synth LSD & Real LSD & Synth DMR & Real DMR \\
    \midrule
    E3~Nercessian~\cite{nercessian2020} & 5.964 & 6.792 & 0.507 & 0.512 \\
    E4~Pepe~\cite{pepe2020}             & 5.809 & 6.623 & 0.504 & 0.516 \\
    \midrule
    AC1~TCN$+$BiLSTM   & \textbf{0.853} & 1.508 & \textbf{0.947} & \textbf{0.910} \\
    AC2~TCN$+$GRU      & 0.860 & \textbf{1.506} & 0.946 & \textbf{0.910} \\
    AC3~TCN$+$Conformer~\cite{gulati2020} & 0.870 & 1.576 & 0.945 & 0.906 \\
    \midrule
    A0~Proposed $\star$ & 1.442 & 1.941 & 0.928 & 0.885 \\
    A1~w/o Room Awareness  & 2.897 & 4.875 & 0.832 & 0.740 \\
    A2~with Pref Loss  & 1.703 & 2.243 & 0.884 & 0.851 \\
    \bottomrule
  \end{tabular}%
  }
\end{table}

\subsection{Architecture Comparison}
\label{sec:arch}

Fig.~\ref{fig:arch_compare} shows a comprehensive architecture
comparison.
\begin{figure}[tp]
  \centering
  \includegraphics[width=\linewidth]{fig10_arch_compare.pdf}
  \caption{Architecture comparison: LSD, DMR, RTF, and LSD domain gap
           across A0 (proposed) and AC1--AC3 variants.}
  \label{fig:arch_compare}
\end{figure}

The three AC variants converge to comparable in-distribution performance
and outperform the proposed A0 on dual-target LSD (0.853--0.870 vs.\
1.442) and DMR (0.945--0.947 vs.\ 0.928).
This is the expected consequence of the AC variants' dense 128-bin
output head, which is unconstrained and can fit arbitrary
frequency-domain targets, in contrast to A0's 7-band parametric-EQ head, whose output is
restricted to a sum of Gaussian-approximated peaking-EQ magnitude
responses with $f_c$, gain, and $Q$ as the only free parameters per
band.
The two design choices serve different purposes: AC variants
demonstrate the architectural ceiling achievable with the same encoder
backbone and training data, while A0's PEQ head produces parameters that can be mapped to embedded
PEQ control interfaces without post-hoc dense-response fitting or
impulse-response inversion.
AC1 (BiLSTM) is bidirectional and is therefore incompatible with
causal streaming deployment.
AC2 (GRU) is causal and remains within the same real-time RTF range.
AC3 (Conformer) has 418{,}914 parameters versus A0's 203{,}447, more
than 2$\times$ as many.
A0 also shows a smaller synthetic-to-real domain gap than the dense-output
AC variants and the auxiliary-loss A2 variant (Section~\ref{sec:ood}),
suggesting that the \texttt{room\_vec}-only Stage~A conditioning
generalizes well to unseen acoustic environments.
The final predicted response of A0 is also closer to the room-only target
under the Room-LSD diagnostic (LSD$_\text{room}$~4.689~dB) than
AC1--AC3 (5.226--5.238~dB). This does not make A0 the best overall
dual-objective predictor; rather, it indicates that its two-stage
factorization preserves a more explicit room-correction-like component
while the dense-output AC variants use their flexibility to optimize the
final dual target.
The proposed model is therefore interpreted as the compact
deployment-compatible configuration, not as the highest-scoring
in-distribution architecture.

% ═══════════════════════════════════════════════════════════════════════
\section{DISCUSSION}
% ═══════════════════════════════════════════════════════════════════════

\subsection{The Dual-Objective Trade-off}

Fig.~\ref{fig:pareto} shows the LSD--DMR trade-off across all evaluated
models.
\begin{figure}[tp]
  \centering
  \includegraphics[width=\linewidth]{fig11_pareto.pdf}
  \caption{LSD--DMR trade-off. The proposed model ($\star$) occupies a
           compact 7-band PEQ-output trade-off point. The AC variants attain
           lower LSD and higher DMR with dense 128-bin output that is not
           directly mappable to embedded PEQ control interfaces.}
  \label{fig:pareto}
\end{figure}

Results indicate that tonal shape tracking fidelity emerges primarily
from the input conditioning structure (UserPreferenceEncoder) rather
than from auxiliary loss terms.
Removing the preference input collapses DMR from 0.928 to 0.509
(ablation A3), while adding explicit Stage-B preference and direction
losses (ablation A2) does not help: A2 yields LSD~1.703~dB and
DMR~0.884, both worse than the simplified two-term A0 (LSD~1.442~dB,
DMR~0.928); A0 attains 0.26~dB lower LSD and 0.044 higher DMR than A2
($d_z = -0.486$ for LSD, $d_z = +0.404$ for DMR; both $p < 0.001$).
The additive dual-target formulation $T_\text{dual}$ already encodes
the tonal-shape target directly, and the single $\mathcal{L}_\text{final}$
term suffices for shape tracking; the explicit Stage-B losses appear
to over-constrain the parametric-EQ head.
The simplification has two consequences: first, the loss is reduced
from four terms to two ($\mathcal{L}_\text{final}$ and
$\mathcal{L}_\text{room}$), eliminating two hyperparameters; second,
the proposed A0 attains marginally better in-distribution metrics and
a smaller synthetic-to-real domain gap than A2.
The proposed model nonetheless trails the dense-output AC variants by
$\sim$0.6~dB on dual-target LSD; this is the explicit cost of
constraining the output to 7 PEQ parameters that are directly
mappable to embedded PEQ control interfaces.


\subsection{Real-Time Feasibility}
\label{sec:rtf}

In a practical deployment, the 128-bin room response can be obtained by
playing a calibration sweep through the target space, recording the
microphone output, windowing the resulting RIR, and interpolating the
log-magnitude response to 128~bins. This calibration is an operating
scenario rather than a component timed in the reported RTF measurements,
and it need not be repeated unless the listening position changes
significantly, as room acoustics are quasi-static on practical timescales.

The model was measured on a standard CPU (AMD Ryzen~9800X3D,
single-threaded). The RTF values reported in this paper are generated by
the evaluation script and are computed as inference-only forward latency
divided by the 4-second segment duration. The system therefore requires a
short audio-context buffer for feature extraction and is segment-based
rather than sample-by-sample streaming.
RTF in Table~\ref{tab:main} (denoted $\dagger$) excludes the time required
to accumulate the audio context, extract acoustic features, and estimate
the room-response input.
Since room response changes slowly relative to musical content,
per-segment EQ updates on the order of seconds are acoustically appropriate.
Model inference runs in under 10~ms on CPU, and the compact parameter count
is compatible with future ARM-based embedded deployment studies.

\subsection{Limitations}

Four limitations and one scope boundary merit note.
(i)~Training targets domestic listening environments (room size
3.5--8.0~m, RT60~0.20--0.75~s); performance degrades in larger spaces
such as lecture or concert halls, as the BUT ReverbDB and OpenAIR
out-of-distribution experiments indicate.
(ii)~The system addresses frequency-response coloration only; it does not
attempt to fully resolve late reverberation, modal ringing, or other
time-domain artifacts that cannot be robustly handled by a static
low-order PEQ alone.
(iii)~Paired statistical tests operate at sample level ($N=3{,}000$),
where samples may share source tracks or room configurations, introducing
within-group correlation; effect sizes (Cohen's $d_z$) remain informative
but $p$-values should be interpreted cautiously. Track-level aggregate
analysis ($N_\text{groups}=1{,}306$) confirms the sample-level findings
(Section~\ref{sec:stats}).
(iv)~The tonal shape presets are illustrative examples and do not
represent validated listener preferences; no perceptual listening
evaluation was conducted, and whether objective shape-tracking
improvements translate to perceptual outcomes is left for future work.

% ═══════════════════════════════════════════════════════════════════════
\section{CONCLUSION}
% ═══════════════════════════════════════════════════════════════════════

This report described a feed-forward neural PEQ that jointly performs
room compensation and operator-specified tonal shape tracking in a
single network with Gaussian-approximated PEQ-parameter output,
achieving DMR~0.928, LSD~1.442~dB, and RTF~$<$~0.001 on a held-out test set of 3{,}000
synthetic-RIR samples.
Ablations identify the UserPreferenceEncoder input conditioning as the
primary driver of tonal-shape fidelity ($d_z = +2.376$, $p < 0.001$ vs.\
removing the preference input) and room-awareness cues as essential
for out-of-distribution robustness (real-RIR LSD~4.875~dB without them
vs.\ 1.941~dB with them).
Architectural variants with a dense 128-bin output (BiLSTM, GRU,
Conformer) attain approximately 0.6~dB lower in-distribution LSD but
do not produce parametric PEQ coefficients and are not directly
mappable to embedded PEQ control interfaces; the proposed 7-band PEQ
output is positioned as the deployment-compatible configuration with
the reported performance and a smaller synthetic-to-real domain gap than
the dense-output AC variants and the auxiliary-loss A2 variant.
A more elaborate four-term loss with explicit Stage-B supervision was
investigated and is reported as ablation A2; it does not improve
test-set metrics over the simplified two-term loss used in the proposed
model.
The implementation reported here integrates both objectives in one
feed-forward pass and emits PEQ parameters that can be mapped to embedded
control interfaces without post-hoc dense-response fitting.
The contribution is a documented academic implementation and evaluation
pipeline for an application-driven design space.
Listener-based perceptual evaluation is outside the present scope.


\subsection*{Future Work}

Future directions include: (i)~subjective perceptual evaluation
(MUSHRA, ITU-R~BS.1534) to relate the objective shape-tracking metrics
to listener-rated outcomes; (ii)~integration with a de-reverberation
front-end for environments where late reverberation dominates;
(iii)~a frequency-dependent conflict-weighting strategy to
mitigate gain conflicts at acoustic nulls; and (iv)~end-to-end
training with a time-varying EQ head as in Diff-DEQ~\cite{sarkar2025},
to capture room-induced spectral variation that evolves over time.



\section*{Code Availability}
Implementation scripts, configuration files, and generated result tables
will be made publicly available upon acceptance, subject to third-party
dataset licensing constraints.


% (Author Contributions section omitted — single-author work)

% ── References ───────────────────────────────────────────────────────
\begin{thebibliography}{99}

\bibitem{keidser2006}
G.~Keidser and H.~Dillon,
``What's New in Prescriptive Fittings Down Under?''
in \emph{Hearing Care for Adults}, pp.~133--142, Phonak AG (2006).

\bibitem{nielsen2013}
J.~B.~B.~Nielsen, J.~Nielsen, B.~S.~Jensen, and J.~Larsen,
``Hearing Aid Personalization,''
in \emph{Proceedings of the NIPS Workshop on Machine Learning and
Interpretation in Neuroimaging} (2013).

\bibitem{mondol2022}
S.~I.~M.~M.~Raton~Mondol, H.~J.~Kim, K.~S.~Kim, and S.~Lee,
``Machine Learning-Based Hearing Aid Fitting Personalization Using
Clinical Fitting Data,''
\emph{J.\ Healthc.\ Eng.}, vol.~2022, art.~1667672 (2022).
\url{https://doi.org/10.1155/2022/1667672}

\bibitem{miyoshi1988}
M.~Miyoshi and Y.~Kaneda,
``Inverse Filtering of Room Acoustics,''
\emph{IEEE Trans.\ Acoust.\ Speech Signal Process.}, vol.~36, no.~2,
pp.~145--152 (1988).
\url{https://doi.org/10.1109/29.1509}

\bibitem{pepe2020}
G.~Pepe, L.~Gabrielli, S.~Squartini, and L.~Cattani,
``Designing Audio Equalization Filters by Deep Neural Networks,''
\emph{Appl.\ Sci.}, vol.~10, no.~7, p.~2483 (2020).
\url{https://doi.org/10.3390/app10072483}

\bibitem{pepe2022}
G.~Pepe, L.~Gabrielli, C.~Tripodi, S.~Squartini, and N.~Strozzi,
``Deep Optimization of Parametric IIR Filters for Audio Equalization,''
\emph{IEEE/ACM Trans.\ Audio Speech Lang.\ Process.}, vol.~30,
pp.~1136--1149 (2022).
\url{https://doi.org/10.1109/TASLP.2022.3155289}

\bibitem{cartwright2013}
M.~Cartwright and B.~Pardo,
``Social-EQ: Crowdsourcing an Equalization Descriptor Map,''
in \emph{Proceedings of the International Society for Music Information
Retrieval Conference (ISMIR)} (Curitiba, Brazil) (2013 Nov.).

\bibitem{seetharaman2016}
P.~Seetharaman and B.~Pardo,
``Audealize: Crowdsourced Audio Production Tools,''
\emph{J.\ Audio Eng.\ Soc.}, vol.~64, no.~9, pp.~683--695 (2016 Sep.).
\url{https://doi.org/10.17743/jaes.2016.0037}

\bibitem{chu2025}
A.~Chu, P.~O'Reilly, J.~Barnett, and B.~Pardo,
``Text2FX: Harnessing CLAP Embeddings for Text-Guided Audio Effects,''
in \emph{Proceedings of the IEEE International Conference on Acoustics,
Speech and Signal Processing (ICASSP)} (Hyderabad, India) (2025).
\url{https://doi.org/10.1109/ICASSP49660.2025.10890334}

\bibitem{stylianou2026}
I.~Stylianou, J.~Francombe, P.~Mart\'inez-Nuevo, S.~E.~Shepstone,
and Z.-H.~Tan,
``Population-Aligned Audio Reproduction with LLM-Based Equalizers,''
\emph{arXiv preprint arXiv:2601.09448} (2026).
\url{https://arxiv.org/abs/2601.09448}

\bibitem{nercessian2020}
S.~Nercessian,
``Neural Parametric Equalizer Matching Using Differentiable Biquads,''
in \emph{Proceedings of the International Conference on Digital Audio Effects
(DAFx)}, pp.~265--272 (Vienna, Austria) (2020 Sep.).

\bibitem{nercessian2021_hc}
S.~Nercessian, A.~Sarroff, and K.~J.~Werner,
``Lightweight and Interpretable Neural Modeling of an Audio Distortion
Effect Using Hyperconditioned Differentiable Biquads,''
in \emph{Proceedings of the IEEE International Conference on Acoustics,
Speech and Signal Processing (ICASSP)}, pp.~890--894 (2021).
\url{https://doi.org/10.1109/ICASSP39728.2021.9413996}

\bibitem{colonel2022}
J.~T.~Colonel, C.~J.~Steinmetz, M.~Michelen, and J.~D.~Reiss,
``Direct Design of Biquad Filter Cascades with Deep Learning by Sampling
Random Polynomials,''
in \emph{Proceedings of the IEEE International Conference on Acoustics,
Speech and Signal Processing (ICASSP)}, pp.~3104--3108 (2022).
\url{https://doi.org/10.1109/ICASSP43922.2022.9747660}

\bibitem{steinmetz2022}
C.~J.~Steinmetz, N.~J.~Bryan, and J.~D.~Reiss,
``Style Transfer of Audio Effects with Differentiable Signal Processing,''
\emph{J.\ Audio Eng.\ Soc.}, vol.~70, no.~9, pp.~708--721 (2022).
\url{https://doi.org/10.17743/jaes.2022.0025}

\bibitem{szoke2019}
I.~Sz\H{o}ke, M.~Sk\'{a}cel, L.~Mo\v{s}ner, J.~Paliesek, and J.~\v{C}ernock\'{y},
``Building and Evaluation of a Real Room Impulse Response Dataset,''
\emph{IEEE J.\ Sel.\ Top.\ Signal Process.}, vol.~13, no.~4, pp.~863--876 (2019).
\url{https://doi.org/10.1109/JSTSP.2019.2917582}

\bibitem{openair2010}
S.~Shelley and D.~T.~Murphy,
``OpenAIR: An Interactive Auralization Web Resource and Database,''
in \emph{Proceedings of the 129th Audio Engineering Society Convention},
paper~8226 (San Francisco, USA) (2010 Nov.).

\bibitem{sarkar2025}
P.~Sarkar and P.~Lindborg,
``Diff-DEQ: Differentiable Dynamic Equalization for Studio-Quality Speech
Processing,''
in \emph{Proceedings of the European Signal Processing Conference (EUSIPCO)},
pp.~511--515 (2025).
\url{https://doi.org/10.23919/EUSIPCO63237.2025.11226258}

\bibitem{mockenhaupt2024}
F.~Mockenhaupt, J.~S.~Rieber, and S.~Nercessian,
``Automatic Equalization for Individual Instrument Tracks Using Convolutional
Neural Networks,''
in \emph{Proceedings of the International Conference on Digital Audio Effects
(DAFx)} (Guildford, United Kingdom) (2024).

\bibitem{perez2018}
E.~Perez, F.~Strub, H.~de~Vries, V.~Dumoulin, and A.~Courville,
``FiLM: Visual Reasoning with a General Conditioning Layer,''
in \emph{Proceedings of the AAAI Conference on Artificial Intelligence},
pp.~3942--3951 (2018).
\url{https://doi.org/10.1609/aaai.v32i1.11671}

\bibitem{defferrard2017}
M.~Defferrard, K.~Benzi, P.~Vandergheynst, and X.~Bresson,
``FMA: A Dataset for Music Analysis,''
in \emph{Proceedings of the International Society for Music Information
Retrieval Conference (ISMIR)}, pp.~316--323 (Suzhou, China) (2017 Oct.).

\bibitem{scheibler2018}
R.~Scheibler, E.~Bezzam, and I.~Dokmanić,
``Pyroomacoustics: A Python Package for Audio Room Simulation and Array
Processing Algorithms,''
in \emph{Proceedings of the IEEE International Conference on Acoustics,
Speech and Signal Processing (ICASSP)}, pp.~351--355 (2018).
\url{https://doi.org/10.1109/ICASSP.2018.8461310}

\bibitem{gulati2020}
A.~Gulati, J.~Qin, C.-C.~Chiu, N.~Parmar, Y.~Zhang, J.~Yu, W.~Han,
S.~Wang, Z.~Zhang, Y.~Wu, and R.~Pang,
``Conformer: Convolution-Augmented Transformer for Speech Recognition,''
in \emph{Proceedings of Interspeech}, pp.~5036--5040 (2020 Oct.).
\url{https://doi.org/10.21437/Interspeech.2020-3015}

\end{thebibliography}

\biography{Yongho Ahn}{photo_yongho.jpg}{%
Yongho Ahn is an undergraduate researcher in the Department of Software
Engineering at Soongsil University, Seoul, Republic of Korea.
His research interests include applied deep learning, edge AI systems,
embedded systems, and AI-based signal processing, with recent work on
neural audio equalization, real-time inference pipelines on embedded system, hardware
prototyping, and Android application development.}

\end{document}